\begin{proof}
\emph{Setting.} By Theorem \claimref{synthesis.signed-directions-and-the-integer-clifford-algebra}(ii), $e_{S}e_{T} = (-1)^{b(S,T)}e_{T}e_{S}$ and $e_{S}^{2} = (-1)^{q(S)}$, so the $h_{S}$ of a family commute, square to $1$, and $\rho = 2^{-m}\prod_{r}(1 + s_{r}h_{S_{r}}) = 2^{-m}\sum_{S \in L}\pm h_{S}$ is an idempotent with $\mathrm{tr}\rho = 1$. \emph{Partial trace.} The commutant of $A$ is generated by $\omega$ and the elements $e_{\{1,\dots,2a\}}e_{j}$, $j > 2a$. A blade $e_{S}e_{T}$, $S \subseteq \{1,\dots,2a\}$, $T \subseteq \{2a+1,\dots,2m\}$, is up to sign $e_{S}e_{\{1,\dots,2a\}}^{|T|}$ times a monomial of the commutant in the $|T|$ generators indexed by $T$, and the normalised trace of the commutant vanishes on that monomial unless $T = \emptyset$. So the trace-preserving conditional expectation onto $A$ keeps the blades with $T = \emptyset$ and kills the others; $\rho_{A} = d_{B}$ times it, that is $\rho_{A} = d_{A}^{-1}\sum_{S \in L_{A}}\pm h_{S}$, and since $L_{A}$ is a subspace on which the signs form a character, $\rho_{A}^{2} = (|L_{A}|/d_{A})\rho_{A}$ and $\mathrm{tr}\rho_{A}^{n} = (|L_{A}|/d_{A})^{n-1}$: flat of rank $d_{A}/|L_{A}|$. \emph{Quarter turns.} For $u^{2} = -1$, $(1 + u)h(1 - u)/2$ is $h$ if $u$ and $h$ commute and $uh$ otherwise, which is $\pm h_{x + w}$: the map $x \mapsto x + b(x,w)w$ on strings, with a change of signs, so states go to states and the moments, which do not depend on the signs, are constant on sign classes.

(i) Let $V_{A}$ be the strings inside $\{1,\dots,2a\}$; $b$ restricted to it is the form of $2a$ directions, nondegenerate (the stored clause (iii), even case), so the string space is the orthogonal sum of $V_{A}$ and $V_{B} = V_{A}^{\perp}$, of dimensions $2a$ and $2b'$ ($V_{B}$ is the set of strings of the commutant's monomials of the lemma above, not the strings inside the last $2b'$ directions). Put $I_{A} = L \cap V_{A} = L_{A}$ and $I_{B} = L \cap V_{B}$, of dimensions $k_{A}$, $k_{B}$. Then $L \subseteq (I_{A} \oplus I_{B})^{\perp}$ and $L/(I_{A} \oplus I_{B})$ is a family of $(I_{A}^{\perp}/I_{A}) \oplus (I_{B}^{\perp}/I_{B})$, of dimensions $2(a - k_{A})$ and $2(b' - k_{B})$, meeting neither summand. A maximal subspace on which the form vanishes and which meets neither summand projects injectively to each, so $a - k_{A} = b' - k_{B} = e$ and it is the graph of a linear bijection $\theta$ with $b(\theta x, \theta y) = b(x, y)$ (the form vanishes on the graph and the field has characteristic two); conversely every such graph is one. There are $|Sp(2e, 2)|$ such $\theta$, $N(a, a - e)$ choices of $I_{A}$ and $N(b', b' - e)$ of $I_{B}$, and $L$ is recovered from $(I_{A}, I_{B}, \theta)$. The count $N(c, k)$: an ordered basis of a $k$-dimensional subspace on which $b$ vanishes is chosen in $\prod_{r=0}^{k-1}(2^{2c - r} - 2^{r})$ ways, and each subspace has $\prod_{r=0}^{k-1}(2^{k} - 2^{r})$ ordered bases; the quotient is $\binom{c}{k}_{2}\prod_{r=0}^{k-1}(2^{c-r} + 1)$. The order of $Sp(2e, 2)$ is counted the same way by symplectic bases. At $k = m$ the total $\prod_{r=1}^{m}(2^{r} + 1)$.

(ii) The number of families containing a given nonzero string $x$ is the number of families of $x^{\perp}/x$, $\prod_{r=1}^{m-1}(2^{r} + 1)$, so a nonzero string lies in a fraction $1/(D + 1)$ of the families, and two distinct commuting nonzero strings lie together in a fraction $1/((D + 1)(D/2 + 1))$. Hence $\langle|L_{A}|\rangle = 1 + (d_{A}^{2} - 1)/(D + 1)$, and dividing by $d_{A}$ gives $(d_{A} + d_{B})/(D + 1)$. For the third moment, $|L_{A}|^{2}$ counts ordered pairs $(x, y)$ of $L_{A}$: one pair $(0,0)$; $3(d_{A}^{2} - 1)$ pairs spanning a line; and $(d_{A}^{2} - 1)(d_{A}^{2}/2 - 2)$ ordered pairs of distinct commuting nonzero strings of $V_{A}$. So $\langle|L_{A}|^{2}\rangle(D + 1)(D + 2) = (D + 1)(D + 2) + 3(d_{A}^{2} - 1)(D + 2) + (d_{A}^{2} - 1)(d_{A}^{2} - 4) = d_{A}^{2}(d_{A}^{2} + d_{B}^{2} + 3D + 1)$, and dividing by $d_{A}^{2}$ gives the stated value. Both values are the blade-state case, with $W = V_{A}$, of clause (iii) of Theorem \claimref{synthesis.quarter-turn-rotors-of-the-even-subalgebra-form-a-state-design-of-order-three}; the count above is an independent derivation of that case. At $n = 2$ the value is the one Fact \claimref{synthesis.imported-haar-moments-of-random-pure-state} states for the mean purity. The same Fact at $n = 3$: the identity gives $d_{A}d_{B}^{3}$, the cycle $\eta$ gives $d_{A}^{3}d_{B}$, its inverse $d_{A}d_{B}$, and the three transpositions $d_{A}^{2}d_{B}^{2}$ each, over $D(D + 1)(D + 2)$: the same.

(iii) The values are evaluations of (i). For the limit put $k = c - e$. With $\binom{c}{k}_{2} = 2^{k(c-k)}\phi_{c}/(\phi_{k}\phi_{c-k})$, $\prod_{r<k}(2^{c-r} + 1) = 2^{kc - k(k-1)/2}\prod_{r<k}(1 + 2^{r-c})$, $|Sp(2e,2)| = 2^{2e^{2} + e}\prod_{r \le e}(1 - 4^{-r})$ and $\prod_{r \le 2c}(2^{r} + 1) = 2^{c(2c+1)}\prod_{r \le 2c}(1 + 2^{-r})$, the power of two in the fraction of families with $e = c - k$ is $2k(c - k) + 2kc - k(k - 1) + 2(c - k)^{2} + (c - k) - c(2c + 1) = -k^{2}$, and the remaining factors tend, as $c \to \infty$ at fixed $k$, to $\phi_{k}^{-2}\prod_{r}(1 - 4^{-r})/\prod_{r}(1 + 2^{-r}) = \phi/\phi_{k}^{2}$. They are bounded uniformly in $c$ by a constant times $\phi_{k}^{-2} \le \phi^{-2}$, so $d^{n-1}\langle\mathrm{tr}\rho_{A}^{n}\rangle_{\mathrm{fin}} = \sum_{k}p_{k}(c)2^{(n-1)k}$ converges by domination to the stated series. \emph{The limit is $G_{n-1}$.} $|L_{A}|^{n-1}$ counts the $(n-1)$-tuples of elements of $L_{A}$; sort them by the subspace $U \subseteq L_{A}$ they span, of dimension $r \le n - 1$. A subspace $U$ of $V_{A}$ lies in some family only if $b$ vanishes on it, and then in the families of $U^{\perp}/U$, a fraction $\prod_{j=m-r+1}^{m}(2^{j} + 1)^{-1}$ of all; there are $N(a, r)$ such $U$; and the tuples spanning a given $U$ are the $r \times (n-1)$ matrices of rank $r$, $\prod_{j<r}(2^{n-1} - 2^{j})$ of them. This gives the displayed sum (at $a = b' = 2$, $n = 4$: $1 + \tfrac{105}{17} + \tfrac{70}{17} = \tfrac{192}{17} = 64 \cdot \tfrac{3}{17}$). At $a = c$, $m = 2c$: $N(c, r) = \binom{c}{r}_{2}\prod_{j<r}(2^{c-j} + 1) = 2^{r(c-r)}2^{rc - r(r-1)/2}(1 + o(1))/\phi_{r}$ and $\prod_{j=2c-r+1}^{2c}(2^{j} + 1) = 2^{2cr - r(r-1)/2}(1 + o(1))$, so their quotient tends to $2^{-r^{2}}/\phi_{r} = 2^{-r(r-1)/2}/\prod_{j=1}^{r}(2^{j} - 1)$; times $\prod_{j<r}(2^{n-1} - 2^{j}) = 2^{r(r-1)/2}\prod_{j<r}(2^{n-1-j} - 1)$ this is $\prod_{j<r}(2^{n-1-j} - 1)/\prod_{j=1}^{r}(2^{j} - 1) = \binom{n-1}{r}_{2}$. The sum has at most $n$ terms, so the limit is $\sum_{r}\binom{n-1}{r}_{2} = G_{n-1}$. The check compares the displayed sum with (i) at small sizes and the values at $c = 16, 24$ with $G_{n-1}$.

(iv) The parity sector is the $+1$ eigenspace of $h_{v}$ on the $D$-dimensional module; in the occupation basis of the lemma's tensor factors it is spanned by the products of an even with an even and an odd with an odd state. The Schur argument of Fact \claimref{synthesis.imported-haar-moments-of-random-pure-state} applies to the sector as a space of dimension $D/2$: the average of $n$ copies is the symmetriser of the sector over $\prod_{l<n}(D/2 + l)$, and $\mathrm{tr}\rho_{A}^{n}$ is the trace against the cyclic permutation of the $A$ factors. For $\sigma \in S_{n}$ the trace of $P_{\sigma}$ composed with that permutation over $n$ copies of the sector counts the basis labels $(\alpha_{r}, \beta_{r})$ with $\alpha$ constant on the cycles of $\sigma^{-1}\eta$, $\beta$ constant on those of $\sigma$ and $\alpha_{r}$, $\beta_{r}$ of equal parity for each $r$; the parity is then constant on the orbits of the group generated by $\sigma$ and $\sigma^{-1}\eta$, which contains $\eta$ and is transitive, so there are $2$ choices of parity and $(d_{A}/2)^{\#(\sigma^{-1}\eta)}(d_{B}/2)^{\#(\sigma)}$ labels for each. At $n = 2$: $2[(d_{A}/2)^{2}(d_{B}/2) + (d_{A}/2)(d_{B}/2)^{2}]/((D/2)(D/2 + 1)) = (d_{A} + d_{B})/(D + 2)$. For the sector average: the families containing $v$ are the families of $v^{\perp}/v$, $\prod_{r=1}^{m-1}(2^{r} + 1)$ of them; a string $x \notin \{0, v\}$ lies in such a family only if it is even, and then in a fraction $1/(D/2 + 1)$ of them; $V_{A}$ has $d_{A}^{2}/2 - 1$ nonzero even strings and does not contain $v$. So the mean of $|L_{A}|/d_{A}$ is $(1 + (d_{A}^{2}/2 - 1)/(D/2 + 1))/d_{A} = (d_{A} + d_{B})/(D + 2)$. For the third moment: two distinct nonzero even strings $x, y$ of $V_{A}$ with $b(x, y) = 0$ span with $v$ a three-dimensional subspace on which $b$ vanishes ($v \notin V_{A}$), and lie together in a fraction $1/((D/2 + 1)(D/4 + 1))$ of the sector families. On the $E = d_{A}^{2}/2$ even strings of $V_{A}$, $b(x, y) = |x \cap y|$ has radical spanned by the full string $v_{A}$ of $A$, so the number of ordered pairs of distinct commuting nonzero even strings is $(E - 2) + (E - 2)(E/2 - 2) = (E - 2)(E/2 - 1)$. With $\alpha = d_{A}/2$, $\beta = d_{B}/2$, $\Delta = D/2 = 2\alpha\beta$: $\langle|L_{A}|^{2}\rangle(\Delta + 1)(\Delta + 2) = (\Delta + 1)(\Delta + 2) + 3(2\alpha^{2} - 1)(\Delta + 2) + 4(\alpha^{2} - 1)^{2} = \Delta^{2} + 6\alpha^{2}\Delta + 4\alpha^{2} + 4\alpha^{4}$, and dividing by $d_{A}^{2} = 4\alpha^{2}$ gives $(\alpha^{2} + \beta^{2} + 3\alpha\beta + 1)/((\Delta + 1)(\Delta + 2))$, which is the sector's Haar value $2[\alpha\beta^{3} + \alpha^{3}\beta + 3\alpha^{2}\beta^{2} + \alpha\beta]/(\Delta(\Delta + 1)(\Delta + 2))$ and equals $(d_{A}^{2} + d_{B}^{2} + 3D + 4)/((D + 2)(D + 4))$. The sign of $h_{v}$ is fixed on the orbit because even quarter turns commute with $e_{v}$.

(v) The quarter turn about $e_{r}e_{j}$ sends a string to its image under the transposition $(r\,j)$ (it adds $\{r,j\}$ exactly when the string contains one of $r$, $j$), so bivector turns act through the symmetric group on the directions and the orbit of the family of the matching $\{1,2\},\dots$ is the set of matching families, $(2m - 1)!! = \prod_{r=1}^{m-1}(2r + 1)$. Since $2r + 1 = 2^{r} + 1$ for $r = 1, 2$ and $2r + 1 < 2^{r} + 1$ for $r \ge 3$, this equals $\prod_{r=1}^{m-1}(2^{r} + 1)$ exactly for $m \le 3$. Permutations preserve the size of a string, and the symmetric group is transitive on the strings of each size, so the orbits on the classes $\{x, x + v\}$ are exactly the size classes $\{2k, 2m - 2k\}$, $1 \le k \le \lfloor m/2\rfloor$: $\{2, 6\}$ and $\{4\}$ at $m = 4$, $\{2, 8\}$ and $\{4, 6\}$ at $m = 5$. For the last statement: $(1 + e_{w})(1 + e_{w})^{\sim} = 1 + e_{w}e_{w}^{\sim} + e_{w} + e_{w}^{\sim}$ is $2$ exactly when $e_{w}^{\sim} = -e_{w}$, and for these generators $e_{w}^{\sim} = (-1)^{q(w)}e_{w}$. For an even string, $q(x + v) = q(x) + q(v) + b(x, v) = q(x) + m \pmod 2$, since $q(v) = m(2m - 1)$ and $b(x, v) = |x|$: for $m$ odd one of $x$, $x + v$ has $q = 1$; for $m$ even $q$ is a function on classes, and a transvection about $w$ with $q(w) = 1$ preserves it, $q(x + w) = q(x) + q(w) + b(x, w) = q(x)$ when $b(x,w) = 1$. For even $m \ge 4$ both values occur (sizes $2$ and $4$).

(vi) A flat state of rank $2^{e}$ has entropy $e\ln 2$; the means follow from (i), and $\ln 2 = \sum_{k \ge 1}1/(k2^{k})$ lies between its twelfth partial sum and that sum plus $1/(13 \cdot 2^{12})$, inside $(\tfrac23, \tfrac{7}{10})$. Page's values are the formula of Fact \claimref{singularity.imported-page-average-entropy-subsystem}.

\emph{Comparison with the stored results.} Theorem \claimref{synthesis.eviction-keeps-one-replica-saddle}(iii), (iv) uses the mean purity only, which (ii) reproduces; Theorem \claimref{synthesis.eviction-keeps-one-of-catalan-many-replica-saddles}(iii), (iv) uses the moments at every index and their Catalan limit, which (iii) does not reproduce from $n = 4$; Theorem \claimref{synthesis.eviction-deficit-against-page-deficit} uses Page's average, which (vi) does not reproduce. The sector form of the crossing value follows from (iv) at $d_{A} = d_{B} = d$: $2d/(d^{2} + 2) = (2/d)/(1 + 2d^{-2})$.
\end{proof}
